% ============================================================
\section{Concerns (CRN)}\label{sec:concerns}
% ============================================================

\subsection{Usuarios y clientes}

\begin{tabla}{|L{1.1cm}|L{5.15cm}|L{2.15cm}|Y|}
\hline
\filacab \cab{ID} & \cab{Concern} & \cab{STK} & \cab{Consecuencia arquitectónica} \\ \hline
\endhead
CRN-01 &
Cuando intento una jugada que el juego no me deja hacer, quiero entender al momento por qué no puedo. Si tengo que ponerme a leer un párrafo para saberlo, me aburro y lo dejo. &
STK-01, STK-09 &
Exige retroalimentación inmediata (previsualización de barreras, resaltado de casillas legales) y obliga a validar reglas también en el cliente. \\ \hline
CRN-02 &
Cuando coloco una barrera quiero que el juego me conteste al instante. Si tarda más de un segundo en decirme si vale, la partida se me hace pesada. &
STK-01 &
El presupuesto se reparte entre la búsqueda de camino, el viaje de red y la escritura en base de datos. Condiciona qué se persiste durante la partida y qué al final. \\ \hline
CRN-03 &
Quiero que el juego me siga interesando cuando ya sé jugar, pero no quiero que aprender a jugar me cueste más de unos minutos. &
STK-01, STK-06 &
Favorece reglas parametrizables (tamaño de tablero, número de barreras) en lugar de variantes codificadas. \\ \hline
CRN-04 &
No voy a dejar que mi hijo juegue si en el chat puede encontrarse insultos o que un desconocido intente hablarle. &
STK-02, STK-17, STK-19 &
El filtrado debe ejecutarse en el servidor, ya que un cliente modificado podría evadirlo. \\ \hline
CRN-05 &
No quiero dar de mi hijo más datos de los imprescindibles. Si para entrar hay que registrarse y verificar identidad, quiero saber exactamente qué se guarda y para qué. &
STK-02, STK-17 &
Puede requerir consentimiento del tutor o cuentas con datos mínimos, lo que afecta el esquema de cuentas. \\ \hline
CRN-06 &
Si se me corta la conexión a mitad de partida quiero saber qué pasa: si puedo volver, si pierdo automáticamente, o si dejo al otro esperando. &
STK-01, STK-14 &
Exige estado autoritativo persistido, tiempos límite de turno y una política de reconexión o rendición. \\ \hline
\end{tabla}

% ------------------------------------------------------------
\subsection{Negocio y dirección}

\begin{tabla}{|L{1.1cm}|L{5.15cm}|L{2.15cm}|Y|}
\hline
\filacab \cab{ID} & \cab{Concern} & \cab{STK} & \cab{Consecuencia arquitectónica} \\ \hline
\endhead
CRN-07 &
Necesito saber si todo lo que se ha comprometido cabe de verdad en las 14 semanas del curso, y qué se recorta primero si no cabe. &
STK-03, STK-04 &
Obliga a aislar las funcionalidades adicionales tras interfaces, para poder recortarlas sin comprometer el núcleo ni el modo en línea. \\ \hline
CRN-08 &
Cada semana quiero ver algo que funcione y que yo pueda probar, no una explicación de lo que se está construyendo por dentro. &
STK-03, STK-16 &
Favorece construir por rebanadas verticales completas en lugar de capa por capa. \\ \hline
CRN-09 &
Necesito poder justificar cada decisión estructural que tomo. Si dentro de dos meses alguien me pregunta por qué manda el servidor y no el cliente, tengo que poder señalar dónde lo decidí y con qué motivo. &
STK-05, STK-16 &
Requiere mantener la trazabilidad como parte del entregable y no como un anexo posterior. \\ \hline
\end{tabla}

% ------------------------------------------------------------
\subsection{Desarrollo y arquitectura}

\begin{tabla}{|L{1.1cm}|L{5.15cm}|L{2.15cm}|Y|}
\hline
\filacab \cab{ID} & \cab{Concern} & \cab{STK} & \cab{Consecuencia arquitectónica} \\ \hline
\endhead
CRN-10 &
Tengo que decidir quién manda sobre el estado de la partida, el cliente o el servidor. De esa decisión cuelga casi todo lo demás, y cambiarla más adelante me sale carísimo. &
STK-05, STK-15 &
Decisión estructural mayor. El servidor autoritativo protege contra trampas pero añade latencia. \\ \hline
CRN-11 &
Sin WebSockets tengo que definir yo el protocolo sobre TCP: dónde empieza y termina cada mensaje, cómo lo serializo, y qué hago cuando el cliente y el servidor no tienen la misma versión. &
STK-05, STK-07 &
Sin WebSockets no existe delimitación automática de mensajes; hay que resolver la fragmentación y la compatibilidad entre versiones del cliente. \\ \hline
CRN-12 &
Si el cliente y el servidor aplican reglas distintas, el jugador verá válida una jugada que el servidor le rechaza. Necesito que los dos usen exactamente la misma lógica. &
STK-05, STK-10 &
Empuja hacia un único motor de reglas compartido como biblioteca, consumido por ambos lados. \\ \hline
CRN-13 &
Mientras equilibro el juego voy a cambiar el tamaño del tablero y el número de barreras decenas de veces. No puedo depender de que alguien recompile cada vez que quiero probar un valor. &
STK-06 &
Los parámetros de reglas deben vivir en configuración externa y no como constantes en el código. \\ \hline
CRN-14 &
Todavía no he decidido si habrá IA ni conexión LAN, ni cómo se va a llamar el juego. Necesito que el equipo pueda avanzar sin que mi decisión, cuando llegue, les obligue a rehacer lo ya hecho. &
STK-04, STK-07 &
Requiere abstraer el concepto de oponente y el extremo de red, y aislar marca y arte del código. \\ \hline
CRN-15 &
Diseño las pantallas en español, pero el mismo botón en inglés puede ocupar la mitad o el doble. Si maqueto con anchos fijos, se me descuadra la interfaz en cuanto se traduce. &
STK-08, STK-11 &
Prohíbe anchos fijos y texto incrustado en imágenes. Debe decidirse antes de maquetar. \\ \hline
CRN-16 &
Necesito poder traducirlo todo sin tocar el código: textos, fechas, números y formatos. Lo que esté escrito dentro del programa o incrustado en una imagen, yo no lo puedo cambiar. &
STK-11 &
Cadenas externalizadas en catálogos de recursos y Unicode de extremo a extremo, incluido el esquema en SQL Server. \\ \hline
CRN-17 &
Necesito ejecutar y repetir partidas enteras sin pasar por la interfaz. Si para reproducir un fallo tengo que ir haciendo clic en el tablero, no me da tiempo a probar lo suficiente. &
STK-10 &
Exige separar el motor de reglas de la interfaz y que la partida sea determinista y registrable jugada por jugada. \\ \hline
CRN-26 &
El mecanismo del segundo factor todavía no está decidido. Necesito dejarlo detrás de una interfaz para poder cambiarlo después sin tener que tocar el flujo de inicio de sesión. &
STK-05, STK-15 &
Empuja a abstraer la verificación del segundo factor tras una interfaz, para poder cambiar de mecanismo sin tocar el flujo de sesión. \\ \hline
\end{tabla}

% ------------------------------------------------------------
\subsection{Operación y soporte}

\begin{tabla}{|L{1.1cm}|L{5.15cm}|L{2.15cm}|Y|}
\hline
\filacab \cab{ID} & \cab{Concern} & \cab{STK} & \cab{Consecuencia arquitectónica} \\ \hline
\endhead
CRN-18 &
Me preocupa que alguien me reporte solo por haberle ganado y acabe sancionándome sin que nadie haya mirado lo que ocurrió de verdad. &
STK-01, STK-15 &
Si solo se cuentan reportes, hace falta un límite por usuario o una ponderación por reputación. Si se quiere revisión real, obliga a persistir contexto de chat y a construir un módulo administrativo. \\ \hline
CRN-19 &
Si el servidor se cae, tengo que poder recuperar las cuentas, el progreso y las partidas. Necesito que se defina qué información es aceptable perder, porque guardarlo todo al instante cuesta tiempo de respuesta. &
STK-13 &
Transacciones, respaldos y una política explícita de qué información es aceptable perder ante una caída. \\ \hline
CRN-20 &
Si mañana se decide añadir LAN, no quiero tener que rehacer el núcleo. El extremo de red tiene que poder configurarse desde fuera, no estar fijo en el código. &
STK-05, STK-14 &
Configuración externa del extremo de red y del modo de despliegue, en lugar de valores fijos en el código. \\ \hline
CRN-21 &
Si el correo o el servidor externo dejan de responder, necesito saber qué hace el sistema: reintentar, dejar la cuenta pendiente o degradarse de forma controlada. Lo que no puede es quedarse colgado. &
STK-13, STK-14 &
Obliga a manejar fallos de terceros mediante reintentos, estados de cuenta pendiente o degradación controlada del servicio. \\ \hline
\end{tabla}

% ------------------------------------------------------------
\subsection{Seguridad, auditoría y regulación}

\begin{tabla}{|L{1.1cm}|L{5.15cm}|L{2.15cm}|Y|}
\hline
\filacab \cab{ID} & \cab{Concern} & \cab{STK} & \cab{Consecuencia arquitectónica} \\ \hline
\endhead
CRN-22 &
El transporte es TCP en claro. Si no ciframos, cualquiera en la misma red puede leer las credenciales y secuestrar la sesión de un jugador. &
STK-15 &
Exige decidir el cifrado sobre el socket y el almacenamiento seguro de contraseñas y secretos del segundo factor. \\ \hline
CRN-23 &
Doy por hecho que alguien va a modificar el cliente, así que voy a probar con uno alterado. Necesito comprobar que no puede hacer jugadas ilegales ni ver información oculta del rival. &
STK-10, STK-15 &
Refuerza el servidor autoritativo: el cliente propone y el servidor decide. \\ \hline
CRN-24 &
Necesito saber hasta dónde puede parecerse el juego a Quoridor sin infringir derechos de terceros. Si me equivoco, habrá que cambiar nombre, arte o redacción de reglas con el proyecto ya avanzado. &
STK-17, STK-18 &
Restringe nombre, arte y redacción de reglas, y obliga a mantenerlos desacoplados para poder cambiarlos tarde. \\ \hline
CRN-25 &
Antes de aceptar el juego en nuestro catálogo exigimos clasificación por edad, un empaquetado conforme a nuestras normas y una política de chat que cumpla la regulación de cada región. &
STK-19 &
Puede imponer requisitos de empaquetado, control parental o la desactivación del chat en ciertas regiones. \\ \hline
\end{tabla}
